% SPDX-License-Identifier: CC-BY-4.0
% Copyright (C) 2026 Fabian Januszewski, Abhijit (the Langlandscape team)
% S01-T1: What JL says.  source: code/results/s01/brandt_w2.json
% math: mathematics/ label ex:s01-masschecks ; stage tex: plan/staged/stage-01-quaternion-foundations.tex, thm:s01-JLclassical (cited)
\begin{frame}{What Jacquet--Langlands says: the case $D=11$}
  % source: code/results/s01/brandt_w2.json  (data.T1.D = 11)
  \begin{columns}[T]
    \begin{column}{0.48\textwidth}
      \textbf{Quaternion side.} $B_{11}$ definite, ramified at $\{11,\infty\}$.\\[2pt]
      $X_B(\Ohat^\times)$ is a finite set:
      \begin{itemize}
        \item % data.T1.h = 2
          $h=2$ points,
        \item % data.T1.gammas = [4, 6]
          automorphism groups of order $4$ and $6$,
        \item % data.T1.mass = "5/12" = data.T1.mass_expected; ex:s01-masschecks
          mass $\tfrac14+\tfrac16=\tfrac5{12}=\tfrac{11-1}{24}$.
      \end{itemize}
      \medskip
      \textbf{Modular side.} $X_0(11)$ has genus % data.T1.X0_11.genus = 1
      $1$ (Weil curve 11a1).
    \end{column}
    \begin{column}{0.48\textwidth}
      \textbf{Brandt matrix} % data.T1.T2 = [[1,2],[3,0]]
      \[
        T_2=\begin{pmatrix}1&2\\3&0\end{pmatrix},\qquad
        \text{eigenvalues } 3,\ -2 .
      \]
      % data.T1.T2_eigenvalues = [3, -2]; data.T1.X0_11.a2_mf = -2; data.T1.X0_11.a2_11a1_ellap = -2
      {\footnotesize\begin{tabular}{@{}lcc@{}}
        \toprule
        & $T_2$, $B_{11}$ & $a_2$, $X_0(11)$\\
        \midrule
        Eisenstein (const.) & $3$ & $3$\\
        newform 11a & $-2$ & $-2$\\
        \bottomrule
      \end{tabular}}
    \end{column}
  \end{columns}
  \vspace{4pt}
  \begin{block}{No map of spaces}
    The correspondence is between \emph{Hecke eigensystems}, not a map between the
    two spaces: $\dim=2$ on one side, genus $1$ on the other.
  \end{block}
  {\tiny Status: \emph{computed} (Brandt $T_2$, Weil curve 11a1); the JL correspondence itself is \emph{cited}.}
\end{frame}
